\pagenumbering{roman}
\setcounter{page}{1}

\selecthungarian

%----------------------------------------------------------------------------
% Abstract in Hungarian
%----------------------------------------------------------------------------
\chapter*{Kivonat}\addcontentsline{toc}{chapter}{Kivonat}

Számos gráfadathalmazban a csúcsokhoz szöveg tartozik, például egy hivatkozási hálózatban a cikkek absztraktjai. Amikor egy ilyen gráfhoz új csúcs kerül, a meglévő csúcsokkal való kapcsolatai már ismertek lehetnek, mielőtt a saját szövege elérhetővé válna. Dolgozatomban azt vizsgálom, hogy egy nyelvi modell képes-e egy új csúcs szövegét a szomszédságából származó információk alapján generálni. Ezt az információt soft promptok, vagyis a modellnek átadott folytonos vektorok közvetítik. A feladatot egy hivatkozási gráfon vizsgálom: egy cikk absztraktját az általa hivatkozott cikkek alapján generálom, miközben a cikk saját címe és absztraktja rejtve marad.

Az ogbn-arxiv hivatkozási gráfot a Semantic Scholar hivatkozási adataival bővítettem, és minden cikk címéből és absztraktjából embeddinget számítottam. Egy gráfokon működő neurális hálózat (GNN) a hivatkozott cikkek embeddingjeiből becsüli meg a célcikk embeddingjét. Egy prefix adapter ezt a becslést key és value vektorokká alakítja a Qwen3-14B minden attention rétege számára. A nyelvi modellt low-rank adaptation (LoRA) segítségével finomhangoltam az absztraktok generálására. A GNN-alapú becslőt és a szöveggeneráló komponenseket először külön tanítottam, majd együtt finomhangoltam.

Az értékelés során egy nyelvi modell alapú bíráló (LLM judge) méri a generált absztraktok tartalmi egyezését az eredeti absztraktokkal, és értékeli a szöveg minőségét. Az azonosítási teszt azt vizsgálja, hogy a generált szöveg alapján azonosítható-e a célcikk a vizsgált cikkek közül. A tanult becslést a hivatkozott cikkek embeddingjeinek átlagával, valamint diagnosztikai kontrollokkal hasonlítom össze, hogy megvizsgáljam a szomszédsági információ szerepét a generálásban.

A generált absztraktok általában koherensek, és a sikeres esetekben tükrözik a célcikk kutatási kontextusát, valamint a célcikkre jellemző kifejezéseket tartalmaznak. A tanult GNN-alapú becslő a hivatkozott cikkek embeddingjeinek átlagához képest javítja a tartalmi egyezést. A generált absztrakt akkor is tűnhet hihetőnek az adott szakterületen, ha nem adja vissza hűen a célcikk módszereit vagy eredményeit.


\vfill
\selectenglish


%----------------------------------------------------------------------------
% Abstract in English
%----------------------------------------------------------------------------
\chapter*{Abstract}\addcontentsline{toc}{chapter}{Abstract}

In many graph datasets, nodes are associated with text, such as the abstracts of the papers in a citation network. When a new node is added to such a graph, its connections to existing nodes may be known before its own text is available. This thesis investigates whether soft prompts, continuous vectors supplied to a language model, can convey enough neighbourhood information to generate a new node's text. I study this task on a citation graph: the abstract of a paper is generated from the papers it cites, while its own title and abstract are withheld.

I enriched the ogbn-arxiv citation graph with citation records from Semantic Scholar and represented the title and abstract of each paper by a text embedding. A graph neural network (GNN) estimates the embedding of the target paper from the embeddings of its references. A prefix adapter turns this estimate into key and value vectors for every attention layer of Qwen3-14B, which is adapted with low-rank adaptation (LoRA) and generates the abstract. The GNN estimator and the text-generation components are first trained separately and then fine-tuned jointly.

Evaluation uses a language-model judge to assess content agreement with the original abstracts and to rate writing quality. An identification test checks whether the generated texts distinguish their target papers. Comparisons with reference averaging and diagnostic input controls examine how neighbourhood conditioning affects generation.

The generated abstracts are generally coherent and, in successful cases, reflect the target paper's research context and contain distinguishing vocabulary. The learned GNN estimator improves content agreement over averaging the reference embeddings. A generated abstract may, however, describe a plausible study in the relevant area without faithfully reconstructing the target paper's methods or findings.


\vfill
\selectthesislanguage

\newcounter{romanPage}
\setcounter{romanPage}{\value{page}}
\stepcounter{romanPage}